

\medskip
\noindent\SGAInlineRefTarget{sga3:iii:lemma:0-1-11}\textbf{Lemma 0.1.11.}
\textit{Let \(X,X'\) be two \(S\)-schemes. There is a commutative
diagram}

\SGAThreeNativeDiagram[0.58\textwidth]
  {native_diagrams/exp3/diagram_015_exp3_p012_lemma0111_product_square.tex}
  {Compatibility of \(L_X\) and \(X^+\) with products.}
  {fig:exp3-lemma-0111-product}
  {Polo--Gille re-edition, Expos\'e III}

\textit{Proof.}\footnote{Editor's note: the proof has been expanded.}
For every \(S\)-scheme \(Y\),
\[
\begin{aligned}
 \operatorname{Hom}_S(Y,X^+\times_S X'^+)
 &=
 \operatorname{Hom}_S(Y,X^+)
 \times\operatorname{Hom}_S(Y,X'^+)\\
 &\simeq
 \operatorname{Hom}_S(Y_{\mathscr J},X)
 \times\operatorname{Hom}_S(Y_{\mathscr J},X')\\
 &=
 \operatorname{Hom}_S\bigl(Y,(X\times_S X')^+\bigr).
\end{aligned}
\]
Thus
\[
 X^+\times_S X'^+\simeq(X\times_S X')^+.
\]

Now let \(Y\) be a scheme over \(X^+\times_S X'^+\), represented by a
morphism
\[
 h:Y_{\mathscr J}\longrightarrow X\times_S X'.
\]
Let \(p,q\) be the projections from \(X\times_S X'\) to \(X,X'\), and
put \(f=p\circ h\), \(g=q\circ h\). Since
\[
 \Omega^1_{(X_0\times_{S_0}X'_0)/S_0}
 \simeq
 p_0^*(\Omega^1_{X_0/S_0})
 \oplus q_0^*(\Omega^1_{X'_0/S_0})
\]
(EGA~IV\textsubscript{4}, 16.4.23), there is a natural isomorphism
\[
\begin{aligned}
 &\operatorname{Hom}_{\mathscr O_{Y_0}}
   \bigl(f_0^*(\Omega^1_{X_0/S_0}),\mathscr J\mathscr O_Y\bigr)
 \times
 \operatorname{Hom}_{\mathscr O_{Y_0}}
   \bigl(g_0^*(\Omega^1_{X'_0/S_0}),\mathscr J\mathscr O_Y\bigr)\\
 &\qquad\simeq
 \operatorname{Hom}_{\mathscr O_{Y_0}}
 \bigl(h_0^*(\Omega^1_{(X_0\times_{S_0}X'_0)/S_0}),
       \mathscr J\mathscr O_Y\bigr).
\end{aligned}
\]
In other words,
\[
 L_X(Y)\times L_{X'}(Y)\simeq L_{X\times_S X'}(Y).
\]

\medskip
\noindent\SGAInlineRefTarget{sga3:iii:remark:0-1-12}\textbf{Remark 0.1.12.}\footnote{Editor's note: this remark
and the corollary that follows have been added.}
Let \(\mathscr C\) be a category stable under fiber products, let \(S\)
be an object of \(\mathscr C\), and let \(T_1,T_2\) be two objects over
\(S\). For \(i=1,2\), let \(L_i\) and \(X_i\) be two objects over
\(T_i\), as in the diagram below.

\SGAThreeNativeDiagram[0.52\textwidth]
  {native_diagrams/exp3/diagram_016_exp3_p012_remark0112_fiber_objects.tex}
  {The objects \(L_i,X_i\) over \(T_i\), and \(T_i\) over \(S\).}
  {fig:exp3-remark-0112-objects}
  {Polo--Gille re-edition, Expos\'e III}

There is then a natural isomorphism
\[
 (L_1\times_{T_1}X_1)\times_S(L_2\times_{T_2}X_2)
 \simeq
 (L_1\times_S L_2)
 \times_{T_1\times_S T_2}(X_1\times_S X_2).
\]
Consequently, the preceding lemma gives:

\medskip
\noindent\SGAInlineRefTarget{sga3:iii:corollary:0-1-13}\textbf{Corollary 0.1.13.}
\textit{Let \(X_1,X_2\) be two \(S\)-schemes. There is a commutative
diagram of isomorphisms:}

\SGAThreeNativeDiagram[0.96\textwidth]
  {native_diagrams/exp3/diagram_017_exp3_p013_corollary0113_isomorphisms.tex}
  {The product isomorphisms of \SGARefLink{sga3:iii:corollary:0-1-13}{Corollary 0.1.13}.}
  {fig:exp3-corollary-0113}
  {LNM 151, p.~88}

We can now state:

\medskip
\noindent\SGAInlineRefTarget{sga3:iii:proposition:0-2}\textbf{Proposition 0.2.}
\textit{For every \(S\)-scheme \(X\), one can define a left action of
the \(X^+\)-abelian group \(L_X\) on the \(X^+\)-object \(X\), with
the following properties:}
\begin{enumerate}
\item[\SGAInlineRefTarget{sga3:iii:proposition:0-2:item:i}(i)] \textit{this action makes \(X\) formally principal
homogeneous under \(L_X\) over \(X^+\); that is, the morphism}
\[
 L_X\times_{X^+}X\longrightarrow X\times_{X^+}X
\]
\textit{is an isomorphism of \(X^+\)-functors;}

\item[\SGAInlineRefTarget{sga3:iii:proposition:0-2:item:ii}(ii)] \textit{the action is functorial in the \(S\)-scheme
\(X\): for every \(S\)-morphism \(f:X\to W\), the following diagram
is commutative;}

\SGAThreeNativeDiagram[0.55\textwidth]
  {native_diagrams/exp3/diagram_018_exp3_p013_proposition02_functoriality_square.tex}
  {Functoriality of the action in \SGARefLink{sga3:iii:proposition:0-2:item:ii}{Proposition 0.2(ii)}.}
  {fig:exp3-proposition-02-functoriality}
  {LNM 151, p.~88}

\item[\SGAInlineRefTarget{sga3:iii:proposition:0-2:item:iii}(iii)] \textit{the action commutes with fiber products: for all
\(S\)-schemes \(X_1,X_2\), the following diagram is commutative.}
\end{enumerate}

\SGAThreeNativeDiagram[0.99\textwidth]
  {native_diagrams/exp3/diagram_019_exp3_p013_proposition02_product_diagram.tex}
  {Compatibility of the action with fiber products.}
  {fig:exp3-proposition-02-products}
  {LNM 151, p.~88}

\textit{Proof.}\footnote{Editor's note: the original has been modified
and expanded to account for the preceding additions; in particular,
those additions incorporate, with details, the content of original
page 89.}
Parts \SGARefLink{sga3:iii:proposition:0-2:item:i}{(i)} and \SGARefLink{sga3:iii:proposition:0-2:item:ii}{(ii)} follow respectively from \SGARefLink{sga3:iii:corollary:0-1-9}{Corollaries 0.1.9} and
\SGARefLink{sga3:iii:corollary:0-1-10}{0.1.10}. To prove \SGARefLink{sga3:iii:proposition:0-2:item:iii}{(iii)}, write
\[
 P(X)=L_X\times_{X^+}X
\]
for every \(S\)-scheme \(X\). Applying \SGARefLink{sga3:iii:proposition:0-2:item:ii}{(ii)} to the projections
\[
 p_i:X_1\times_S X_2\longrightarrow X_i
\]
gives the following commutative squares.

\SGAThreeNativeDiagram[0.55\textwidth]
  {native_diagrams/exp3/diagram_020_exp3_p013_proposition02_projection_square.tex}
  {The square induced by each projection \(p_i\).}
  {fig:exp3-proposition-02-projection-square}
  {LNM 151, p.~88}

For \(i=1,2\), these yield the commutative square

\SGAThreeNativeDiagram[0.55\textwidth]
  {native_diagrams/exp3/diagram_021_exp3_p014_proposition02_final_square.tex}
  {The product square in the proof of \SGARefLink{sga3:iii:proposition:0-2:item:iii}{Proposition 0.2(iii)}.}
  {fig:exp3-proposition-02-final-square}
  {LNM 151, pp.~90--91}

Combining this square with \SGARefLink{sga3:iii:corollary:0-1-13}{Corollary 0.1.13} shows that its vertical
arrow is an isomorphism and gives precisely the commutative diagram
asserted in \SGARefLink{sga3:iii:proposition:0-2:item:iii}{(iii)}.

